\documentclass[11pt]{amsart}
\usepackage{amsmath,amssymb,amsthm,microtype}
\usepackage[hidelinks]{hyperref}
\numberwithin{equation}{section}
\newtheorem{theorem}{Theorem}[section]
\newtheorem{lemma}[theorem]{Lemma}
\newtheorem{proposition}[theorem]{Proposition}
\theoremstyle{definition}
\newtheorem{definition}[theorem]{Definition}
\newtheorem{construction}[theorem]{Construction}
\title{Golden-ratio equilibria for two facilities with atomic clients}
\date{}
\begin{document}
\begin{abstract}
Krogmann, Lenzner, Skopalik, Uetz, and Vos introduced two-stage facility location with atomic clients and left open whether the golden ratio guarantees an approximate subgame-perfect equilibrium. We settle this question for two facilities with the same finite location catalog. For arbitrary positive client weights, we construct a selection of exact client Nash equilibria supporting a $\phi$-approximate facility equilibrium, where $\phi=(1+\sqrt5)/2$. For rational input given by explicit coverage sets, the construction runs in polynomial time and returns rational client strategies. The approximation factor matches the known two-facility lower bound. The proof uses finite menus of client equilibria and an overlap estimate for two selected locations. We show that these local choices exclude every closed response cycle that could obstruct the required equilibrium.
\end{abstract}
\maketitle

\section{Introduction}\label{sec:introduction}
A facility in a two-stage location game chooses its location before clients decide which facility to use. Its payoff therefore depends on the client equilibrium selected after the move. Atomic clients may have several Nash equilibria at a given placement, yielding different facility loads. To establish subgame-perfect equilibrium, we must choose a facility placement together with client equilibria at all possible placements.

Krogmann, Lenzner, Skopalik, Uetz, and Vos~\cite{krogmann2024} introduced this atomic model. They established exact existence for equal client weights and a golden-ratio lower bound for weighted clients, leaving the corresponding existence guarantee open. We prove that the golden ratio suffices when two facilities share the same finite location set. For rational input, we also construct the facility placement and its client continuations in polynomial time. Thus the two-facility common-catalog case attains the known lower bound.

\subsection{The model}
We use the atomic client model of~\cite[Section~1.1]{krogmann2024}, with two facilities and a common location catalog. Describing eligibility by sets permits the location and client sets to be distinct.
\begin{definition}\label{def:model}
An instance consists of a finite client set $G$, positive weights $w_g$, a finite nonempty location set $S$, and coverage sets $C_s\subseteq G$ for $s\in S$. Both labeled facilities choose a location in $S$, simultaneously and with co-location permitted. After a placement $(s,t)$, each covered client chooses one of its accessible facilities, independently of the other clients. An uncovered client is not served. A client minimizes its expected realized facility load, including its own weight; a facility maximizes its expected attracted weight.

A complete continuation rule $\sigma$ assigns an exact client Nash equilibrium to each ordered placement. For $\alpha\ge1$, a placement $(s,t)$ and such a rule form an $\alpha$-approximate subgame-perfect equilibrium if
\begin{equation}\label{eq:spe}
 L_1(\sigma(r,t))\le\alpha L_1(\sigma(s,t))\quad(r\ne s),\qquad
 L_2(\sigma(s,r))\le\alpha L_2(\sigma(s,t))\quad(r\ne t).
\end{equation}
\end{definition}
The requirement is the existence of one complete rule. It does not impose stability under every choice of client equilibria. The multiplicative inequalities in~\eqref{eq:spe} also cover a zero on-path payoff.

Write $w(X)=\sum_{g\in X}w_g$ and $R_s=w(C_s)$. At a placement $(s,t)$, let $A=w(C_s\setminus C_t)$ and $B=w(C_t\setminus C_s)$. Enumerate the common clients by $1,\ldots,k$, put $W=\sum_iw_i$, and let $p_i$ be client $i$'s probability of using the first facility. The two expected loads and their difference are
\begin{equation}\label{eq:loads}
 L_s=A+\sum_iw_ip_i,\qquad L_t=B+\sum_iw_i(1-p_i),\qquad\Delta=L_s-L_t.
\end{equation}
The difference between client $i$'s conditional costs is $\Delta+w_i(1-2p_i)$. Consequently, the exact client Nash conditions are
\begin{equation}\label{eq:ne}
 p_i=1\Rightarrow\Delta\le w_i,\quad
 p_i=0\Rightarrow\Delta\ge-w_i,\quad
 0<p_i<1\Rightarrow\Delta=w_i(2p_i-1).
\end{equation}
These are individual mixed strategies, with no division or correlation of a client's atomic weight.

\subsection{The proof idea and its relation to earlier work}
The argument has two main ingredients. First, we construct finite menus of exact client equilibria, using the same menus to support a placement and to evaluate every deviation from it. Minimizing the deviator's load within these menus defines a response map on locations. Second, we analyze a closed response cycle under the assumption that its menus contain no supporting placement. The overlap inequalities require the pairs of sufficiently heavy clients to form a star. Two selected locations then bound the common client weight along the cycle, excluding both a closed path of decreasing reaches and a return to a small deviation threshold.

For pure client assignments, we use the repair operation of Vos~\cite[Lemma~29]{vos2023}. Our overlap argument builds on the analysis in~\cite[Section~8.3]{vos2023}; the proof below treats every cycle length and supplies the required mixed client equilibria. The squared-load potential establishes termination of local client repairs. To obtain the polynomial-time facility construction, we additionally bound the number of exchanges in the mixed-equilibrium algorithm.

\noindent\textbf{Organization.}
Section~\ref{sec:proof} defines the four equilibrium families, states the structural lemma, and deduces the main theorem. Appendix~\ref{app:cycle} proves the structural lemma by contradiction. Appendices~\ref{app:cross-chord} and~\ref{app:polynomial-chord} prove the local mixed-equilibrium result and its polynomial construction.

\noindent\textbf{Conventions and notation.}
Throughout, $\phi=(1+\sqrt5)/2$, $q=\phi^{-1}$, $a=q^2=1-q$, and $c=\phi/2$. Client weights are positive real numbers in the existence statements and positive binary rationals in the complexity statements. All client randomizations are independent, all placement profiles are labeled, and all tie choices in a construction are fixed by index.

\section{The golden-ratio theorem}\label{sec:proof}
We first construct the client equilibria needed at a pair of locations. We then show that the resulting finite menus contain a placement satisfying both facilities' deviation bounds. The main theorem follows by choosing the client continuations from these menus.

\subsection{Local equilibria and repair}\label{pure-repair}
A pure improvement moves weight $w$ from a facility of load $H$ to one of load $L<H$ precisely when $w<H-L$. Both resulting loads lie strictly between $L$ and $H$. This elementary fact gives the following form of the repair from~\cite[Lemma~29]{vos2023}.
\begin{lemma}\label{lem:repair}
Suppose a pure client assignment has loads $L<H$, and each common client initially on the lower-load side has weight at least $H-L$. Repeatedly move a largest strictly improving client from the original higher-load side to the original lower-load side, stopping when no such client remains or when the load order first reverses. The output is an exact pure client Nash equilibrium. Both final loads are at least $L$, and the original lower-load side finishes strictly below $H$. At most $k$ clients move.
\end{lemma}
\begin{proof}
Before a reversal, the load gap decreases, and the moved weights are nonincreasing: a client too heavy to improve cannot become improving as that gap decreases. If the order reverses after moving weight $w$, the new gap is less than $w$. Every previously moved client and every common client initially on the lower side has weight at least $w$, so none can move back. The remaining clients are on the new lower side. If no reversal occurs, the stopping condition and the initial weight bound give the same conclusion. Each move leaves both loads in the preceding load interval and each client moves at most once.
\end{proof}
In particular, putting every common client on one side always gives either an equilibrium or a seed to which Lemma~\ref{lem:repair} applies. Thus every placement has a pure default continuation.

For comparable reaches and sufficiently small common clients, the next lemma gives one equilibrium in which both facilities receive a prescribed fraction of the opposite reach. We construct it by assigning an individual mixing probability to each common client.
\begin{lemma}\label{lem:chord}
Consider two locations with reaches $U\ge V>cU$ and common client weight $W<qV$. If every common client has weight at most $aU$, there is an exact independently mixed client Nash equilibrium with
\begin{equation}\label{eq:chord}
 L_U\ge qV,\qquad L_V\ge qU.
\end{equation}
For rational input with $k$ common clients, one such equilibrium can be constructed with $O((k+1)^4)$ rational operations and polynomial intermediate encoding length.
\end{lemma}
Appendix~\ref{app:cross-chord} proves existence and gives all equilibrium coordinates. Appendix~\ref{app:polynomial-chord} bounds the number of exchanges used to obtain them. For example, with private loads $40,39$ and three common clients of weight $20$, all three common clients can use the first facility with probability $39/80$. The resulting loads $(277/4,279/4)$ satisfy~\eqref{eq:chord} for reaches $100,99$. Pure equilibria or equilibria with only two mixing clients do not satisfy these two bounds.

\subsection{Four fixed families of client equilibria}\label{sec:menu}
For each pair of coverage sets, we now construct a finite menu from the original list of weighted clients. The construction precedes the choice of a response map, so its members remain available throughout the cycle argument.
\begin{construction}\label{con:menu}
At a pair let $A,B$ denote exclusive customer loads, let
$w_1,\ldots,w_k$ be the common customer weights, and let $W=\sum_iw_i$.
The reaches are $A+W,B+W$. Fix all tie choices by customer index.
Construct $F(s,t)$ as follows. A label reversal uses the reflection of this same menu.
\begin{description}

\item[\textbf{P.}] \emph{Guarded pure seeds.}
Use the two assignments putting all common customers on one side. Also use,
for every common customer $h$, both assignments putting only $h$ on one side
and all other common customers on the other. Retain a seed if it is already
an NE. Otherwise, let its initial loads be $Q<P$, and retain a repaired
output only if each common customer initially on the lower side has weight
at least $P-Q$. Move a largest strictly improving customer from the original
higher side to the original lower side, stopping at the first reversal of
load order or when no improvement remains. Lemma~\ref{lem:repair} proves
that the output is an NE and preserves the initial lower load as a lower
bound. Unqualified seeds are discarded. The all-common seeds make this
pure submenu nonempty: either the occupied side is already lower, or the
other side has no common customer and the guard is vacuous.
There are at most $2(k+1)$ outputs.

\item[\textbf{E.}] \emph{Equal-reach half splitting.}
If $A=B$, let each common customer independently use each side with
probability $1/2$. Each customer is indifferent. In particular this includes
loads $(R_s/2,R_s/2)$ at every co-location $(s,s)$.

\item[\textbf{T.}] \emph{Two designated mixers.}
Choose distinct common customers $h,j$. Orient the endpoints so that
$A\le B$, and put all other common customers on the $B$ side. Set
$Z=W-w_h-w_j$ and $D=B-A+Z$. If $D\le\min(w_h,w_j)$, give $h,j$
probabilities $(1+D/w_h)/2,(1+D/w_j)/2$ on the $A$ side.
The expected load gap is $D$, so the designated customers are indifferent
and the others prefer the lower expected-load side. This is an exact NE.
If $A=B$, include both orientations. This family has at most
$2\binom{k}{2}$ members, including boundary probabilities.

\item[\textbf{C.}] \emph{A two-location quota equilibrium.}
Order the reaches as $U\ge V$. Whenever
\[
 V>cU,\qquad W<qV,\qquad \max_iw_i\le aU,
\]
include one deterministic output of the strong-chord construction of
Lemma~\ref{lem:chord}.
It gives $L_U\ge qV$ and $L_V\ge qU$ in
$O((k+1)^4)$ rational arithmetic operations. For $k=0$ the maximum-weight
condition is vacuous. At a co-location, retain E for the equal-half
argument; it is not, in general, a witness for the stronger $qR$ quotas.
If the strict chord hypotheses hold for a pair of distinct locations,
include the chord output even when the two reaches are equal. This call
depends only on the pair's actual reaches and common weights, never on
global deviation quotas.

\end{description}

\end{construction}
Every member satisfies~\eqref{eq:ne}. For family T, the expected load difference is $D$, so the two designated clients are indifferent and the other common clients prefer the side with smaller expected load. Lemmas~\ref{lem:repair} and~\ref{lem:chord} justify P and C, respectively. Family E follows directly from~\eqref{eq:ne}. The menu contains $O((k+1)^2)$ members and at least one pure member.

\subsection{The structural lemma and the main theorem}\label{sec:criterion}
For an oriented pair, let $L_s(e)$ denote the load received at $s$ in menu member $e$. Define
\begin{equation}\label{eq:du}
 u(s,t)=\min_{e\in F(s,t)}L_s(e),\qquad
 d(t)=\max_{s\in S}u(s,t),\qquad
 b(t)\in\mathop{\rm argmax}_{s\in S}u(s,t).
\end{equation}
The maximum defining $d(t)$ includes the entire catalog. The finite-family minimum is sufficient for our construction, whether or not smaller payoffs are possible at equilibria outside the family.

The next lemma compares the two loads at a candidate placement with the largest deviation minima. It supplies the placement from which we construct the complete continuation rule.
\begin{lemma}\label{lem:menu-cover}
For every instance of Definition~\ref{def:model}, the menus in Construction~\ref{con:menu} contain a pair $(s,t)$ and a member $e\in F(s,t)$ such that
\begin{equation}\label{eq:quota}
 L_s(e)\ge qd(t),\qquad L_t(e)\ge qd(s).
\end{equation}
In fact, every directed cycle of $b$ contains such a pair of locations.
\end{lemma}
Since $b$ maps the finite nonempty catalog to itself, it has a directed cycle. Fix such a cycle and suppose that its menus contain no pair satisfying~\eqref{eq:quota}. Its maximum $d$-value is positive, since a zero maximum makes any equal-half co-location a suitable witness. After scaling this maximum to one, every response edge has a responder load exceeding $c$ times the source reach and an incumbent load below $q$. Appendix~\ref{app:cycle} first shows that the pairs of clients heavier than $q/2$ form a star. We then choose two consecutive locations and apply inclusion--exclusion to their common clients. This excludes both a closed path of decreasing reaches and an exit from the path. Every equilibrium used in the contradiction belongs to the fixed menus, including the mixed equilibria in family C.

The load bounds in Lemma~\ref{lem:menu-cover} allow us to choose consistent client continuations after every facility deviation.
\begin{theorem}\label{thm:main}
Every instance of Definition~\ref{def:model} has a $\phi$-approximate subgame-perfect equilibrium. With $n$ positive binary-rational client weights and $N$ explicitly listed locations, one can construct such an equilibrium with $O(N^2(n+1)^4)$ rational operations and polynomial intermediate bit length. The output gives rational on-path and unilateral-deviation continuations and a polynomial-time default rule for every other placement.
\end{theorem}
\begin{proof}
Choose $(s,t),e$ by Lemma~\ref{lem:menu-cover}, and after each unilateral deviation select a menu member minimizing the deviator's load. By~\eqref{eq:du} and~\eqref{eq:quota}, the deviating load is at most $\phi$ times the corresponding on-path load. Since the two deviation families consist of disjoint ordered placements, these selections, completed elsewhere by Lemma~\ref{lem:repair}, define a complete continuation rule. The menu bounds below give the stated complexity.
\end{proof}
Krogmann et al.~\cite[Theorem~5]{krogmann2024} give two-facility instances with no $\alpha$-approximate equilibrium for any prescribed $\alpha<\phi$. Thus the existence factor in Theorem~\ref{thm:main} is best possible. Their unrestricted instances lie within the common-catalog class. Their decision hardness result allows varying facility numbers and distinct location restrictions; it is a different statement from the two-facility construction proved here.

\subsection{The constructive bound}\label{sec:algorithm}
Construct all pair menus, store one member attaining each minimum $u(s,t)$, and compute $d$ by~\eqref{eq:du}. Scan the menu members for~\eqref{eq:quota}; Lemma~\ref{lem:menu-cover} guarantees success. Return the selected member and the minimizing continuations for the at most $2N-2$ actual unilateral deviations. At all other placements, use the all-common-client seed and Lemma~\ref{lem:repair}.

At a pair with $k$ common customers P has at most $2(k+1)$ outputs and
costs $O((k+1)^3)$ rational operations: each guarded repair moves at most
$k$ customers and a direct scan per move costs $O(k+1)$.
T has $O((k+1)^2)$ outputs and costs $O((k+1)^3)$ including direct Nash
verification. E is linear. Appendix~\ref{app:polynomial-chord} proves an
$O((k+1)^4)$ bound for C. Therefore all menus, minima, maxima, cycle
selection and scanning require $O(N^2(n+1)^4)$ rational operations.
Explicitly storing every probability vector uses
$O(N^2(n+1)^3)$ rational entries, still polynomial.

Every pure load is a sum of input weights. Every exchange coordinate in
Appendix~\ref{app:polynomial-chord} is an input endpoint or a signed sum
of input weights and the slice target. The later mixer gaps divide such
sums by an integer at most $n$, and probabilities additionally divide by
an input weight. Thus all intermediate rational encodings have length
polynomial in the total binary input length. Comparisons with $q$ are
exact: for rational $z\ge0$, the sign of $z-q$ is the sign of
$z^2+z-1$, and $z<0$ is immediate. Using $q^2=1-q$ reduces all other
irrational thresholds to this form. Zero-payoff cases are handled by
cross-multiplication, without division by zero. This proves the stated bit complexity.

The output is independently checkable using the individual Nash inequalities
in~\eqref{eq:ne} and the actual deviation inequalities~\eqref{eq:spe}. The
verifier need not reconstruct the menus, prove their completeness, or
compute minimum payoffs. The short certificate and the polynomial search
proof are separate parts of the result.


\appendix
\section{Excluding a closed response cycle}\label{app:cycle}
\begin{proof}[Proof of Lemma~\ref{lem:menu-cover}]
We use only members of the menus already constructed, while all deviation maxima continue to range over the full catalog.
\subsection{The normalized no-witness inequalities}
Choose, with arbitrary tie breaking,

\[
b(t)\in\arg\max_s u(s,t)
\]

and let \(T\) be any directed cycle of this map. Suppose no menu witness on a pair of cycle locations satisfies \eqref{eq:quota}. All uses of \(d\) below retain the maximum over the \textbf{full} location set \(S\). The cycle maximum \(M=\max_{s\in T}d(s)\) is positive: if \(M=0\), an equal-half co-location equilibrium at any cycle location already satisfies \eqref{eq:quota}. Divide all weights and loads by \(M\), so

\[
\max_{s\in T}d(s)=1.
\]

\textbf{Scope convention.} Unless explicitly stated otherwise, every location quantified below belongs to \(T\). The value \(d(s)\) continues to maximize over the full set \(S\). In particular, all reach upper bounds and d-value upper bounds below are asserted for cycle locations, not for arbitrary locations outside the cycle. Lemma~\ref{lem:chord} is a separate two-location statement and has its own hypotheses.

Co-location at \(s\), with all covered customers independently choosing each facility with probability \(1/2\), gives load \(R_s/2\) on each facility. Thus failure of \eqref{eq:quota} implies

\begin{equation}
d(s)>cR_s,\qquad R_s<2q\quad(s\in T). \label{eq:old-2-2}
\end{equation}

Let \(R=\max_{s\in T}R_s\). Since \(u(b(s),s)\le R_{b(s)}\) and \(b(T)\subseteq T\),

\begin{equation}
1\le R<2q. \label{eq:old-2-3}
\end{equation}

For every edge \(s\to t=b(s)\) and every member of its fixed menu, write \(P=L_t\), \(Q=L_s\). Then

\begin{equation}
P\ge d(s)>cR_s,\qquad Q<qd(t)\le q. \label{eq:old-2-4}
\end{equation}

Indeed, the responder already exceeds its stability threshold \(qd(s)\); if the incumbent reached \(qd(t)\), \eqref{eq:quota} would hold.


\subsection{Reach bounds and return inequalities}\label{a-useful-bound-on-d}
An all-common-client seed bounds the deviation threshold whenever the source reach is large:


\begin{equation}
R_s\ge q\quad\Longrightarrow\quad d(s)\le R_s. \label{eq:old-3-1}
\end{equation}

To prove this, suppose \(d(s)>R_s\), and put \(t=b(s)\). Let \(E=w(C_t\setminus C_s)\). If \(E\le R_s\), assign all common customers to \(s\). Either this is an equilibrium or pure repair gives an equilibrium with responder load at most \(R_s\). The seed or its guarded output belongs to P, contradicting \(u(t,s)=d(s)>R_s\). Consequently \(E>R_s\). Every common customer then strictly prefers \(s\): its cost there is at most \(R_s\), whereas choosing \(t\) costs at least \(E+w_g\). Thus all actual customer equilibria coincide; the nonempty menu consists of that equilibrium, with responder load \(E=d(s)\) and incumbent load \(R_s\). Equation \eqref{eq:old-2-4} forces \(d(t)>\phi R_s\), and hence \(R_s<q\). This proves \eqref{eq:old-3-1}.

The same argument gives the following return rule: if \(d(s)>R_s\), the edge from \(s\) has a unique customer equilibrium, the incumbent receives \(R_s\), and the next d-value exceeds \(\phi R_s\).


\subsection{Every client visible on the cycle has weight below \texorpdfstring{$q$}{q}}\label{no-visible-customer-has-weight-at-least-q}

Suppose a customer of weight \(H\ge q\) is visible at a location of \(T\). Whenever its current location is \(s\), \eqref{eq:old-2-4} forces it to remain accessible at \(b(s)\), since otherwise the incumbent load would be at least \(H\ge q\). It therefore covers the entire cycle.

For any edge \(s\to t\), let \(S_{st}\) be the total weight of its common customers other than this customer. If

\[
R_t\ge R_s-S_{st},
\]

assign the heavy customer to \(s\) and all other common customers to \(t\). The loads are

\[
Q=R_s-S_{st}\ge H,\qquad P=R_t-H<H\le Q,
\]

where \(R_t<2q\le2H\). The heavy customer is stable because \(Q-H\le P\), equivalently \(Q\le R_t\); the other common customers choose the lower-load side and are stable. This pure equilibrium contradicts \(Q<q\) in \eqref{eq:old-2-4}. Thus every edge satisfies

\[
R_t<R_s-S_{st}\le R_s,
\]

which is impossible on a cycle. Therefore

\begin{equation}
w_g<q\quad\text{for every customer visible on }T. \label{eq:old-4-1}
\end{equation}


\subsection{Heavy pairs: size, avoidance, and a chord lemma}\label{heavy-pairs-size-avoidance-and-a-chord-lemma}

Call a customer \textbf{large} if \(w_g>\tau=q/2\), and put

\[
\mathcal H_s=\{g\in C_s:w_g>\tau\}.
\]


\subsubsection{At most two large customers at a location}\label{at-most-two-large-customers-at-a-location}

In a pure member of the menu of an edge, two large customers belonging to the source cannot both be assigned to the responder. If they were, with source reach \(r\), then \(Q<r-q\). Both transferred common customers have weight at least \(P-Q\), so

\[
2P\le r+Q<2r-q.
\]

Together with \(2P>\phi r\), this gives \((2-\phi)r>q\), hence \(r>\phi\), contradicting \(r<2q<\phi\).

If a source covered three large customers, \(Q<q\) would force at least two of them to the responder. An equal-reach edge cannot avoid this argument by being mixed: its equal-half equilibrium has equal loads, each at least \(d(s)>\phi R_s/2>3/4>q\), again contradicting \eqref{eq:old-2-4}. A pure member of P exists in the unequal-reach case. Consequently

\begin{equation}
|\mathcal H_s|\le2. \label{eq:old-5-1}
\end{equation}


\subsubsection{Avoidance of a large pair}\label{avoidance-of-a-large-pair}

If \(s\) contains two large customers and \(t\) contains neither, then

\begin{equation}
R_t<qd(s)\le q. \label{eq:old-5-2}
\end{equation}

The exclusive weight of \(s\) against \(t\) is \(A>q\). Put all common customers at \(t\). If \(R_t\le A\), this is a pure equilibrium with loads \(A,R_t\). If \(R_t>A\), repair gives both loads at least \(A>q\). Hence \(R_t\ge qd(s)\) would satisfy both thresholds in \eqref{eq:quota}.

In particular, all large pairs that occur on \(T\) intersect pairwise.


\subsubsection{Chord retaining exactly one member of a pair}\label{chord-retaining-exactly-one-member-of-a-pair}

In the normalized no-witness situation, for cycle locations \(s,t\), suppose \(s\) has large customers \(h,j\), \(t\) contains \(h\) but not \(j\), and \(R_s\le R_t\). Let \(H=w_h,J=w_j\). Assign \(h\) to \(s\) and all other common customers to \(t\). If their other-common total is \(\sigma\), the loads are

\[
U=R_s-\sigma\ge H+J>q,\qquad V=R_t-H.
\]

The customer \(h\) is stable since \(U-H\le V\). If \(V\le U\), this is a pure equilibrium. Since \(U>q\ge qd(t)\), the no-witness assumption then forces \(V<qd(s)\), and hence \(d(s)>\phi V=\phi(R_t-H)\). If \(V>U\), then

\[
V-U\le R_t-2H-J<H,
\]

because \(R_t<2q<3H+J\). Repair then gives both loads above \(q\), and therefore a quota witness because \(d(s),d(t)\le1\), a contradiction. Thus the displayed implication is valid under the stated no-witness hypothesis:

\begin{equation}
d(s)>\phi(R_t-H). \label{eq:old-5-3}
\end{equation}

In particular, a pair-source best response cannot have at least the source reach while retaining exactly one member of that pair: the constructed equilibrium has incumbent load above \(q\), while every menu member has responding load at least \(d(s)\).


\subsection{Capacity of a location containing a large pair}\label{capacity-of-a-location-containing-a-large-pair}
We next bound the reach of a location from the weights of its two large clients.


If \(\mathcal H_s=\{h,j\}\), then

\begin{equation}
R_s<q+\max(H,J). \label{eq:old-6-1}
\end{equation}

First let \(r=R_s\ge1\). Equal reaches on its edge are impossible by \eqref{eq:old-2-4}. In a pure menu member, if two or more common customers went forward, then \(2P\le r+Q<r+q\), while \(2P>\phi r\); this would imply \(r<1\). Therefore exactly one common customer goes forward. Since the two large customers cannot both remain at the incumbent, the transferred customer is one of them, say of weight \(H\). Then \(r-H=Q<q\), proving \eqref{eq:old-6-1}.

Now suppose \(r<1\) but \(r\ge q+\max(H,J)\). Then \(r>3q/2\) and \(P>\phi r/2>3/4\). Exactly one of the two large customers goes forward, say \(H\), and the other \(J\) remains. If the responder also covered \(J\), it would have the same pair. If its reach were at least one, the already proved case would give its reach \(<q+\max(H,J)\le r<1\), impossible. If its reach were below one, its load would be \(<1-J<1-q/2<3/4\), also impossible. Thus \(J\) is inaccessible to the responder.

Reassign only \(H\) forward and all other common customers to the incumbent. Let their total be \(\sigma\); the loads are

\[
Q_0=r-H\ge q,\qquad P_0=R_{b(s)}-\sigma,
\]

and

\[
\sigma\le r-H-J,\qquad P_0\ge d(s)-r+H+J.
\]

The original equilibrium condition for \(H\) implies \(P_0\le r\), so \(H\) is still stable. If \(P_0\ge Q_0\), the new profile is a pure equilibrium, contradicting \eqref{eq:old-2-4}. If \(P_0<Q_0\), then

\[
Q_0-P_0<(2-\phi/2)r-2H-J<H,
\]

using \(r<1\) and \(3H+J>2q>2-\phi/2\). Repair produces a responder load strictly below \(Q_0\). But

\[
Q_0=r-H<\phi r/2<d(s),
\]

since \(H>q/2>ar/2\). This contradicts the definition of \(d(s)=u(b(s),s)\), proving \eqref{eq:old-6-1}.


\subsection{No edge connects two locations with the same large pair}\label{no-edge-connects-two-locations-with-the-same-large-pair}

Let two locations have the same pair of weights \(H\ge J>q/2\). Write their private weights as \(A\le B\), the total of other common customers as \(W\), and label their reaches \(r_\ell\le r_h\). Thus

\[
r_\ell=A+H+J+W,\qquad r_h=B+H+J+W.
\]

Set \(D=B-A+W\). Capacity gives

\[
0\le D\le B+W=r_h-H-J<q-J<J\le H.
\]

Assign all common customers other than the large pair to the higher-reach location. Let each large customer of weight \(w\in\{H,J\}\) choose the lower-reach location with probability \((1+D/w)/2\). This is an equilibrium: both large customers are indifferent and all other common customers choose the lower expected-load side. Its loads are

\[
L_h=r_\ell-W-(H+J)/2,\qquad L_\ell=r_h-(H+J)/2.
\]

In either possible orientation as source and responder, the responder load is at most

\[
R_{\rm source}-(H+J)/2<cR_{\rm source}<d({\rm source}).
\]

Here \(H+J>q>aR_{\rm source}\), because \(R_{\rm source}<2q\) and \(2aq<q\). This contradicts a best-response edge. Hence

\begin{equation}
\boxed{\text{No b-edge joins two locations with the same large pair.}} \label{eq:old-7-1}
\end{equation}




\subsection{The large-pair family is a star}\label{the-large-pair-family-is-a-star}

Let \(\mathcal E=\{\mathcal H_s:|\mathcal H_s|=2\}\). Its members pairwise intersect by \eqref{eq:old-5-2}. A pairwise-intersecting family of two-element sets either has a common member or contains exactly the three pair types of a triangle. Suppose such a triangle occurs, on customers \(h,j,k\).

Every location of reach at least \(q\) must meet all three pair types, by \eqref{eq:old-5-2}. By \eqref{eq:old-5-1} its own large set is therefore one of the triangle pairs. Choose a maximum-reach location \(A\), of reach \(R\ge1\), and suppose it has pair \(\{h,j\}\).

We claim that both \(H\) and \(J\) exceed \(a\). Suppose, for example, that \(H\le a\). Choose another triangle-pair location containing \(h\) but not \(j\). Apply Appendix~\ref{chord-retaining-exactly-one-member-of-a-pair} with \(A\) as the higher-reach endpoint. Its load \(R-H\ge1-a=q\), and the other endpoint's load exceeds \(q\); either the displayed profile or repair is stable, a contradiction.

Every pair location \(s\) has \(d(s)>q\). Otherwise \eqref{eq:old-2-2} gives \(R_s<2a\). It cannot have the same pair as \(A\), whose two weights exceed \(a\). If it has another triangle pair, sharing \(g\) with \(A\), then \(w_g<q\), so

\[
R-w_g>1-q=a\ge qd(s).
\]

Appendix~\ref{chord-retaining-exactly-one-member-of-a-pair} again constructs a stable equilibrium, a contradiction.

Thus pair locations form a nonempty b-closed subset: their responders have reach at least their d-value, hence above \(q\), and so must again be pair locations. An edge to a different pair cannot weakly increase reach by Appendix~\ref{chord-retaining-exactly-one-member-of-a-pair}. An edge to the same pair is impossible by \eqref{eq:old-7-1}. Hence every edge in this finite closed subset strictly decreases reach, a contradiction. The triangle is impossible, proving

\begin{equation}
\mathcal E\ne\varnothing\quad\Longrightarrow\quad\bigcap_{E\in\mathcal E}E\ne\varnothing. \label{eq:old-8-1}
\end{equation}


\subsection{Common-h chords, transport, and two-anchor covering}\label{common-h-chords-transport-and-two-anchor-covering}

The rest of the proof repeatedly uses a customer \(h\), of weight \(H\), satisfying

\begin{equation}
H>aR,\qquad K:=R-H<q,\qquad H<q. \label{eq:old-9-1}
\end{equation}

In particular \(R<3H\).


\subsubsection{An upward common-h chord}\label{an-upward-common-h-chord}

For locations \(s,t\) containing \(h\) and \(R_s\le R_t\), assign \(h\) to \(s\) and all other common customers to \(t\). Their loads are

\[
U=R_s-\sigma\ge H,\qquad V=R_t-H.
\]

The customer \(h\) is stable. If \(V\le U\), the profile is an equilibrium. If \(V>U\), then \(V-U\le R-2H<H\), so repair applies and the final load at \(t\) is at most \(V\). Therefore

\begin{equation}
u(t,s)\le R_t-H\le K. \label{eq:old-9-2}
\end{equation}


\subsubsection{A common-h edge with d above q has unique transport}\label{a-common-h-edge-with-d-above-q-has-unique-transport}

Suppose \(s\to t\), both contain \(h\), and \(d(s)>q\). If \(R_t\ge R_s\), the upward common-\(h\) chord gives
\[
d(s)=u(t,s)\le R_t-H\le R-H=K<q,
\]
which contradicts \(d(s)>q\). Hence \(R_t<R_s\). Take a pure menu member from P. If \(h\) stayed at the source, responder load would be at most \(R_t-H\le K<q\), impossible. If another common customer accompanied \(h\) forward, their equilibrium conditions would give

\[
2P\le R_s+Q\le2R_s-H,
\]

so \(P\le R_s-H/2<cR_s<d(s)\), using \(H>aR\ge aR_s\). Thus \(h\) is the only common customer going forward.

Let \(\sigma\) be the other-common total, and \(\Delta=R_s-R_t>0\). The pure loads are \(Q=R_s-H\), \(P=R_t-\sigma\), with \(\sigma<R-q<H\). If \(\sigma\ge\Delta\), the reverse assignment, placing \(h\) at the source and all other common customers at the responder, would be an exact singleton seed in P with responder load \(R_t-H\le K<q\), impossible. Consequently \(\sigma<\Delta\).

By this inequality, \(h\) strictly prefers the responder regardless of other choices: its minimum source cost is \(R_s-\sigma>R_t\). Given this choice, every other common customer strictly prefers the source, whose maximum load is \(R_s-H\le K<q\), below the responder's minimum load \(R_t-\sigma\ge d(s)>q\). Thus \textbf{all}, including mixed, customer equilibria coincide. In particular, the finite-menu minimum equals this unique payoff, and

\begin{equation}
d(s)=R_t-\sigma,\quad Q=R_s-H,\quad 0\le\sigma<R_s-R_t. \label{eq:old-9-3}
\end{equation}


\subsubsection{Two-anchor coverage}\label{general-two-anchor-covering-lemma}
The following inclusion--exclusion estimate prevents a high-reach path from reaching a low deviation threshold.
\begin{lemma}\label{lem:two-anchor}

Suppose an edge \(X\to Y\) has the unique transport \eqref{eq:old-9-3}. Write

\[
x=R_X>y=R_Y,\quad \delta=d(X)=y-\sigma,\quad \varepsilon=d(Y),
\quad B=\phi(y-H).
\]

Assume

\begin{equation}
\delta>B,\quad \varepsilon>B,\quad \delta>\phi H,\quad \varepsilon\le y. \label{eq:old-9-4}
\end{equation}

Then there is no h-containing location \(t\) with

\begin{equation}
B<R_t\le y,\qquad d(t)\le B. \label{eq:old-9-5}
\end{equation}

\end{lemma}
\begin{proof}
For each anchor \(i\in\{X,Y\}\), assign \(h\) to \(t\) and the other common customers to \(i\). If their total \(\sigma_{it}\) satisfied \(\sigma_{it}\le R_t-qd(i)\), the initial loads would be

\[
U_t=R_t-\sigma_{it}\ge qd(i),\qquad V_i=R_i-H\ge y-H=qB\ge qd(t).
\]

The customer \(h\) is stable since \(R_t\le R_i\). If \(V_i\le U_t\), this is an equilibrium. Otherwise repair applies because \(R<3H\); both final loads are at least \(U_t\), which exceeds both required thresholds because \(d(i)>B\ge d(t)\). Thus no solution forces

\[
\sigma_{Xt}>R_t-q\delta,\qquad \sigma_{Yt}>R_t-q\varepsilon.
\]

Their intersection, excluding \(h\), has weight at most \(\sigma\), so

\[
\sigma_{Xt}+\sigma_{Yt}\le R_t-H+\sigma.
\]

It follows that

\[
R_t<L:=q(\delta+\varepsilon)-H+\sigma.
\]

But

\[
B-L=\delta-q\varepsilon+q\sigma-qH
\ge a\delta-qH>0,
\]

where only \(\varepsilon\le y=\delta+\sigma\) is used. Hence \(R_t<B\), contradicting \eqref{eq:old-9-5}. \end{proof}


\subsection{A common heavy client excludes a closed cycle}\label{abstract-center-theorem-and-nondegenerate-stars}
The overlap bound now links a local transport edge to the full cycle. We first prove the conditional statement for one common client and then verify its hypotheses for a nondegenerate star.



\subsubsection{Abstract center theorem}\label{abstract-center-theorem}

Under \eqref{eq:old-9-1}, suppose every location without \(h\) has reach at most \(\phi H\). Then the no-solution closed cycle is impossible.

First, every h-containing location has \(d(s)>H\). Suppose instead \(u\) contains \(h\) and \(d(u)\le H\). Then \(R_u<2qH<\phi H\). For any \(t\) with \(R_t>\phi H\), the assumption makes \(t\) an h-location. If \(d(t)\le\phi H\), put \(h\) at \(u\) and other common customers at \(t\). The loads satisfy \(U\ge H\ge qd(t)\) and \(V=R_t-H>qH\ge qd(u)\). The common-h chord construction, with repair if necessary, gives stability. Thus all locations of reach above \(\phi H\) have d-value above \(\phi H\). This set is nonempty since \(R\ge1>\phi H\), and it is b-closed. It consists entirely of h-locations. On it \eqref{eq:old-9-2} forbids every nondecreasing-reach edge, since \(K<q<\phi H\), a contradiction.

If \(u\to v\) leaves \(h\), \eqref{eq:old-2-4} gives

\[
d(v)>\phi H\ge R_v.
\]

The return rule of Appendix~\ref{a-useful-bound-on-d} makes the next edge's equilibrium unique. Its target \(w\) has reach at least \(d(v)>\phi H\), hence contains \(h\), and

\begin{equation}
d(w)>\phi R_v\ge\phi d(u)>\phi H. \label{eq:old-10-1}
\end{equation}

Choose an h-location \(X\) maximizing d among h-locations, and put \(\delta=d(X)\). We have \(\delta>cR\), since the maximum-reach location contains \(h\), and \(\delta>\phi H\): if there is an exit, use \eqref{eq:old-10-1}; otherwise all locations contain \(h\), so \(\delta=1>\phi H\). The responder \(Y=b(X)\) contains \(h\); otherwise \eqref{eq:old-10-1} would produce an h-location with d exceeding \(\phi\delta>\delta\).

Since \(\delta>\phi H>q\), Appendix~\ref{a-common-h-edge-with-d-above-q-has-unique-transport} applies. Use its notation \(x,y,\sigma,\varepsilon\). Set \(B=\phi(y-H)\). Then \(\varepsilon>\phi(x-H)>B\). Since \(\varepsilon\le\delta\), we also have \(x<H+q\delta\) and hence \(B<\delta\). Moreover \(B>K\): if \(H\le R/2\), use \(y\ge\delta>cR\); if \(H\ge R/2\), use \(y\ge\delta>\phi H\), which gives \(B>H\ge K\). The covering lemma applies because \(y>q\) and hence \(\varepsilon\le y\).

Starting at \(Y\), as long as the path remains among h-locations with d above \(B\), \eqref{eq:old-9-2} forces strictly decreasing reach. A first h-location with d at most \(B\) would still have reach above \(B\) and at most \(y\), contradicting Appendix~\ref{general-two-anchor-covering-lemma}. The path must therefore leave \(h\). At that exit its source has d above \(B\); \eqref{eq:old-10-1} gives a returning h-location with d above \(\phi B\). But

\[
\phi B-\delta
=\phi(\delta-\phi H)+\phi^2\sigma>0,
\]

contradicting the choice of \(\delta\). This proves the common-client argument.


\subsubsection{Nondegenerate stars satisfy its assumptions}\label{nondegenerate-stars-satisfy-its-assumptions}

Suppose \(\mathcal E\) contains at least two distinct pairs, with unique common customer \(h\). Every location of reach at least \(q\) contains \(h\): otherwise it must meet two different pairs by covering their two leaves, producing the triangle already excluded in Appendix~\ref{the-large-pair-family-is-a-star}. Thus every no-h reach is below \(q\).

Apply the maximum-reach first-edge argument of Appendix~\ref{the-maximum-reach-first-edge} below (that argument itself does not assume an empty or single-pair family). It gives a unique forward customer \(g\) on \(A\to b(A)\), of weight exceeding \(R-q\), and other-common weight

\[
\sigma<R-cR=aR/2<aq=q^3<q/2.
\]

Both endpoints have reach above \(q\), so both contain the star center \(h\). If \(g\ne h\), the weight of \(h\), which exceeds \(q/2\), would be included in \(\sigma\), impossible. Thus \(g=h\), and \(H>R-q\ge aR\), while \(H<q\) by Appendix~\ref{no-visible-customer-has-weight-at-least-q}. In particular \(\phi H>q\). All no-h reaches are below \(q<\phi H\), and the common-client argument rules out this case for arbitrary cycle length.


\subsection{The first edge and selection of a single-large-customer anchor}\label{the-first-edge-and-selection-of-a-single-large-customer-anchor}

The first-edge argument below applies to every large-pair family, as used in Appendix~\ref{nondegenerate-stars-satisfy-its-assumptions}. After that argument, only \(\mathcal E=\varnothing\) or \(\mathcal E=\{\{h,j\}\}\) remains after Appendix~\ref{nondegenerate-stars-satisfy-its-assumptions}. Here ``h-only'' means that \(h\) is the location's only \textbf{large} customer; it may have any number of small customers.


\subsubsection{The maximum-reach first edge}\label{the-maximum-reach-first-edge}

Choose \(A_0\) of reach \(R\), and put \(A_1=b(A_0)\). Equation \eqref{eq:old-2-4} gives responder load \(P\ge d(A_0)>cR\) and incumbent load \(Q<q\). An equal-reach responder is impossible because the equal-half equilibrium has \(P=Q\). Thus \(R_1<R\), and take a pure menu member from P.

There must be at least one common customer assigned forward, since otherwise \(Q=R\ge1\). If at least two went forward, their equilibrium conditions would give

\[
2P\le R+Q<R+q,
\]

which, with \(2P>\phi R\), would imply \(R<1\). Therefore exactly one customer \(h\) goes forward. Its weight is

\begin{equation}
H=R-Q>R-q\ge aR>q/2. \label{eq:old-12-1}
\end{equation}

Appendix~\ref{no-visible-customer-has-weight-at-least-q} gives \(H<q\), so \eqref{eq:old-9-1} holds. If \(\sigma_0\) is the other-common weight, then \(P=R_1-\sigma_0>q\), so \(\sigma_0<R-q<H\). As in Appendix~\ref{a-common-h-edge-with-d-above-q-has-unique-transport}, \(\sigma_0\ge R-R_1\) would make the reverse pure assignment an exact singleton seed in P with responder load \(R_1-H<K<q\). Therefore \(\sigma_0<R-R_1\), all actual equilibria coincide by strict dominance, so the menu minimum is their common payoff and

\begin{equation}
d(A_0)=R_1-\sigma_0,\qquad d(A_1)>\phi(R-H)=\phi K. \label{eq:old-12-2}
\end{equation}

If there is a unique large pair, \(h\) belongs to it. Indeed, a maximum-reach location must meet that pair by \eqref{eq:old-5-2}; the presence of a different large customer would create another pair.


\subsubsection{A high h-only anchor always exists}\label{a-high-h-only-anchor-always-exists}

If \(A_0\) is h-only, use it. Otherwise it has the unique pair, so \(A_1\) is h-only by \eqref{eq:old-7-1}, since it contains \(h\). In either case an h-only location \(Z\) has reach

\begin{equation}
R_Z>cR. \label{eq:old-12-3}
\end{equation}


\subsubsection{The case H at least R/2}\label{the-case-h-at-least-r2}

Suppose \(H\ge R/2\). Use the h-only anchor \(Z\) above, with reach \(z\). If any no-h location \(t\) had reach \(r_t>\phi H\), then

\[
r_t>cz,\qquad z>cR\ge cr_t.
\]

Its common weight with \(Z\) satisfies

\[
C\le z-H\le H<qr_t,\qquad C\le z/2<qz.
\]

Every common customer is small, with weight at most \(q/2<a z\), because \(z>cR\ge c\) and \(ac=q/2\). Both reaches exceed \(c^2R>q\). Apply Lemma~\ref{lem:chord} and then \eqref{eq:old-3-1}: the resulting loads meet \(qd(t)\) and \(qd(Z)\), giving stability.

Thus every no-h reach is at most \(\phi H\), and the common-client argument gives a contradiction. This rules out \(H\ge R/2\).


\subsubsection{Selecting consecutive anchors when H is below R/2}\label{selecting-consecutive-anchors-when-h-is-below-r2}

Henceforth \(H<R/2\). By \eqref{eq:old-12-2},

\begin{equation}
d(A_1)>\phi K>cR. \label{eq:old-12-4}
\end{equation}

If \(A_1\) is h-only, select \(X=A_0,Y=A_1\).

The only other possibility is that \(A_0\) is h-only and \(A_1\) has the unique pair. Let \(A_2=b(A_1)\). If \(A_2\) lacked \(h\), then

\[
R_2\ge d(A_1)>\phi K>cR.
\]

The common weight between \(A_0\) and \(A_2\) is at most \(K<qR_2\), and all these common customers are small because \(A_0\) is h-only. Their weights are at most \(q/2<aR\). The strong cross-chord lemma gives stability, using \eqref{eq:old-3-1} at both endpoints, a contradiction. Therefore \(A_2\) contains \(h\). Equation \eqref{eq:old-7-1} makes it h-only. Since \(d(A_1)>q\), Appendix~\ref{a-common-h-edge-with-d-above-q-has-unique-transport} gives unique h-transport along \(A_1\to A_2\). Select \(X=A_1,Y=A_2\).

In either selection, write

\[
x=R_X>y=R_Y>cR,\quad \delta=d(X)=y-\sigma>cR>\phi H,
\quad \varepsilon=d(Y).
\]

The selected edge has unique h-transport, \(0\le\sigma<x-y\), and \(Y\) is h-only. Furthermore

\begin{equation}
y>\phi K. \label{eq:old-12-5}
\end{equation}

For \(Y=A_1\), this follows from \(d(A_1)>\phi K\) and \(d(A_1)\le R_1\). For \(Y=A_2\), it follows from \(R_2\ge d(A_1)>\phi K\). The use of \eqref{eq:old-3-1} in the first case is valid because \(R_1>cR>q\).


\subsection{Final two-anchor closure for the empty and single-pair cases}\label{final-two-anchor-closure-for-the-empty-and-single-pair-cases}

Use \(X,Y,x,y,\delta,\varepsilon,\sigma\) from Appendix~\ref{selecting-consecutive-anchors-when-h-is-below-r2}, and set

\[
B=\phi(y-H).
\]

The edge inequality gives \(\varepsilon>\phi(x-H)>B\). Also \(\varepsilon\le y\) by \eqref{eq:old-3-1}, so \(x<H+qy\). Therefore

\[
\delta=y-\sigma>2y-x>(2-q)y-H>B,
\]

where the last difference is \(q(H-ay)>0\). Finally, \eqref{eq:old-12-5} gives

\begin{equation}
B-K>\phi(R-2H)>0. \label{eq:old-13-1}
\end{equation}

Thus all assumptions \eqref{eq:old-9-4} of the two-anchor covering lemma hold.

If \(H\ge y/2\), any no-h location of reach above \(\phi H\) forms a stable strong cross-chord with the h-only anchor \(Y\), exactly as in Appendix~\ref{the-case-h-at-least-r2}: both reach ratios exceed \(c\), common weight is below \(q\) times either reach, and common weights are at most \(q/2<a y\). Hence all no-h reaches are at most \(\phi H\), and the common-client argument is a contradiction.

It remains that \(H<y/2\). Now

\begin{equation}
B>cy>c^2R>q. \label{eq:old-13-2}
\end{equation}

Any no-h location \(t\) with reach \(r_t>B\) forms a stable strong cross-chord with \(Y\). Indeed,

\[
r_t>cy,\qquad y>cR\ge cr_t,
\]

and, since \(H>aR\ge ay\),

\[
C(Y,t)\le y-H<qy,\qquad C(Y,t)\le y-H<qr_t.
\]

All common customers are small, with weight at most \(q/2<a y\). Both reaches exceed \(q\), so \eqref{eq:old-3-1} upgrades the cross-chord reach thresholds to the exact global d-thresholds. Thus no solution forces

\begin{equation}
h\notin C_t\quad\Longrightarrow\quad R_t\le B. \label{eq:old-13-3}
\end{equation}

Start at \(Y\), which has \(d(Y)>B\), and follow \(b\). As long as a source has d above \(B\), its responder has reach at least that d-value and hence above \(B\). By \eqref{eq:old-13-3} it contains \(h\). An h-to-h edge with d above \(B>K\) strictly decreases reach by \eqref{eq:old-9-2}.

If the path never encountered a d-value at most \(B\), this would give strictly decreasing reaches around a finite closed cycle. Otherwise let \(t\) be the first such location. Its reach is above \(B\), because the preceding source has d above \(B\), and is at most \(y\), by the preceding strict decreases. It contains \(h\), and \(d(t)\le B\). This contradicts the two-anchor covering lemma, Appendix~\ref{general-two-anchor-covering-lemma}.

The empty-pair and single-pair cases are impossible. The nondegenerate-star case was ruled out in Appendix~\ref{abstract-center-theorem-and-nondegenerate-stars}, and Appendix~\ref{the-large-pair-family-is-a-star} excludes all other pair families. Thus the no-witness assumption is false.
\end{proof}


\section{The cross-chord lemma for arbitrarily many common customers}
\label{app:cross-chord}
We prove the existence assertion of Lemma~\ref{lem:chord} by constructing individual equilibrium coordinates. The running time is established in Appendix~\ref{app:polynomial-chord}.
\subsection{Statement and normalization}

Let \(q=(\sqrt5-1)/2\), \(a=q^2=1-q\). Consider two locations with reaches \(R\ge r>\phi R/2\). Their common clients have arbitrary finite positive weights \(w_i\), total \(C\), satisfying

\[
C<qr,\qquad \max_i w_i\le aR.
\]

Each common client chooses a location independently and minimizes expected total weight at its chosen location, including itself. Exclusive clients are forced. Then a possibly mixed Nash equilibrium exists with

\[
L_A\ge qr,\qquad L_B\ge qR.
\]

Here \(A\) is the location with reach \(R\). All probabilities constructed below are individual, independent probabilities; no grouping or correlation of clients is used.

Normalize \(R=1\), write \(r=1-\delta\), so \(0\le\delta<a/2\). Let \(V=2-\delta-C\) be the union weight, and put

\[
U=V-2q=2a-C-\delta,
\qquad Z=V-2qr=2a-C+q^3\delta.
\]

The two lower bounds in~\eqref{eq:chord} are equivalent to the following interval for the expected load difference \(\Delta=L_A-L_B\) is \([-Z,U]\).

Useful consequences are

\[
U>q^4-a\delta>q^4/2>0,\quad Z\ge U,\quad Z>3\delta/2,\quad V/2>q.
\]

For the third inequality, \(Z>q^4+(q+q^3)\delta\), whose difference from \(3\delta/2\) is positive throughout \([0,a/2]\); at the upper endpoint the difference is \(aq^3/4>0\). For the fourth, \(V>1+ar>1+q/2>2q\).


\subsection{Equilibrium coordinates and three mixers}\label{equilibrium-coordinates-and-three-mixers}

Give common client \(i\) coordinate \(c_i=w_i(2p_i(A)-1)\). Its equilibrium conditions are

\[
\begin{array}{ll}
c_i=w_i&\Longrightarrow\Delta\le w_i,\\
c_i=-w_i&\Longrightarrow\Delta\ge-w_i,\\
|c_i|<w_i&\Longrightarrow c_i=\Delta.
\end{array}
\]

The load identity is \(\Delta=\delta+\sum_i c_i\). If there are \(k\ge2\) mixed clients and pure masses \(P_A,P_B\), this gives

\[
(k-1)\Delta=-\delta-P_A+P_B.
\]

Every equilibrium with at least three mixed clients already meets both targets. In the constructions below, a designated mixing probability may equal an endpoint; the same coordinate formulas remain valid by continuity, and such a client may be treated as a designated mixer for the mass inequalities. For \(\Delta=d\ge0\), all mixed weights are at least \(d\), and

\[
C\ge \delta+(2k-1)d+2P_A\ge\delta+5d.
\]

Hence \(d\le(C-\delta)/5<U\). Indeed,

\[
6C+4\delta<6q+(4-6q)\delta<5-2q<10a,
\]

where the last strict inequality is \(q<5/8\). For \(\Delta=-e\le0\), similarly

\[
C\ge(2k-1)e-\delta+2P_B\ge5e-\delta,
\]

so \(e\le(C+\delta)/5<Z\). To verify this last inequality,

\[
6C+(1-5q^3)\delta<6q+(6-16q)\delta\le6q<10a.
\]

These numerical mass implications apply even when the equilibrium in question has only two mixers: whenever \(C\ge\delta+5d\) or \(C\ge5e-\delta\), the same conclusion follows.


\subsection{Largest client and the maximum-square construction}\label{largest-client-and-the-maximum-square-construction}

There is nothing to prove if \(C=0\). Otherwise let \(x\) be a largest common weight and let \(S=C-x\).

If \(\delta>S\), hold \(x\) at \(B\) and choose a pure Nash equilibrium of the remaining clients. Such an equilibrium exists by the usual finite improvement potential. For every intermediate assignment of the remaining clients, the load difference satisfies
\[
L_B-L_A\le x+S-\delta<x,
\]
so client \(x\) remains stable throughout the repair sequence. In particular, the terminal equilibrium has no improvement by \(x\).

Moreover \(L_B\ge r-S>r-\delta=2r-1>q\). If \(L_A<qr\), then \(L_B-L_A>Z>S\), so every remaining shared client must be at \(A\); this gives \(L_A=1-x\ge1-a=q\ge qr\), a contradiction. This settles \(\delta>S\).

Now assume \(\delta\le S\). On the nonempty compact polytope

\[
\mathcal P=\{(c_i)_{i\ne x}: -w_i\le c_i\le w_i,\ \sum_{i\ne x}c_i=-\delta\},
\]

use the polynomial exchange algorithm of Appendix~\ref{app:polynomial-chord} to obtain a vertex at which no feasible two-coordinate endpoint move increases \(\sum_{i\ne x}c_i^2\). A vertex has at most one coordinate strictly between its bounds. If all coordinates are endpoints, assign \(x\) probability \(1/2\); this gives \(\Delta=0\) and is an equilibrium.

Otherwise let the unique interior coordinate have weight \(z\le x\) and value \(d\). Assign \(x\) coordinate \(d\), so \(x\) and \(z\) both mix independently with probability \((1+d/w)/2\). The load identity gives \(\Delta=d\). For a pure \(A\) client of weight \(v\), decreasing its coordinate and increasing the pivot is locally feasible, and pairwise endpoint maximality gives \(d\le v\). For a pure \(B\) client, the opposite perturbation gives \(d\ge-v\). Thus this is a Nash equilibrium.

The same two-coordinate endpoint comparisons give stronger exchange conditions. If \(d>0\), every pure \(A\) client has weight at least \(z\); every pure \(B\) client of weight below \(z\) has weight at most \(d\).

For the first assertion, consider a pure \(A\) client of weight \(v\). If \(2v\ge z-d\), move the pivot to \(z\) and its coordinate from \(v\) to \(v-(z-d)\). The change in the squared objective is

\[
2(z-d)(z-v),
\]

so \(v\ge z\). If \(2v<z-d\), flip its coordinate to \(-v\) and increase the pivot to \(d+2v\). The gain \(4v(d+v)>0\) is impossible.

For a pure \(B\) client of weight \(v<z\), if \(2v\ge z+d\), moving the pivot to \(-z\) and its coordinate to \(-v+(z+d)\) gives positive gain \(2(z+d)(z-v)\), impossible. Otherwise flipping it to \(+v\) and moving the pivot to \(d-2v\) has gain \(4v(v-d)\), forcing \(v\le d\).

For negative gap \(-e\), the symmetric statements hold: every pure \(B\) client has weight at least \(z\); pure \(A\) clients below \(z\) have weight at most \(e\).


\subsection{A bad positive gap cannot persist}\label{a-bad-positive-gap-cannot-persist}

Suppose the constructed equilibrium has gap \(d>U\). If a pure \(A\) client exists, its weight is at least \(z\ge d\). With pure masses \(P_A,P_B\), the two-mixer balance gives \(P_B=\delta+P_A+d\), and hence

\[
C=x+z+2P_A+\delta+d\ge\delta+5d.
\]

Appendix~\ref{equilibrium-coordinates-and-three-mixers} rules this out. Therefore all other clients are pure \(B\), with total weight \(\delta+d\).

Call such a client large if its weight is at least \(z\); every other one has weight at most \(d\). If some latter client has weight \(v\in[d/3,d]\), mix it with \(x,z\). The new gap is \((d-v)/2\), which is within all three weights, and every remaining client stays pure \(B\). This is a three-mixer equilibrium and meets the two load bounds. We can therefore assume every nonlarge client has weight below \(d/3\).

There must be a large client. Otherwise order the small weights decreasingly and take the prefix just before its running sum first reaches \(d\). Let its sum be \(T<d\), and let \(v\) be the next weight. Then \(d-T\le v\), and every weight in the prefix is at least \(v\). Mix this prefix together with \(x,z\), keeping all others at \(B\). If the prefix contains \(m\) clients, the new gap is \((d-T)/(m+1)\), feasible for every mixer. Since each small weight is below \(d/3\), the prefix is nonempty; the resulting equilibrium has at least three mixers and meets the two load bounds.

We next show \(d\le\delta\). If at least two large clients exist, \(\delta+d\ge2z\ge2d\). If exactly one exists, call its weight \(y\ge z\) and let the small total be \(T=\delta+d-y\). If \(\delta+T\le2z\), mix \(x,y,z\) and place every small client at \(A\); the gap is \(-(\delta+T)/2\), feasible for all three mixers, and gives the two load bounds. Otherwise

\[
2\delta+d\ge\delta+T+z>3z\ge3d,
\]

which gives \(\delta>d\).

Change to the pure profile with \(x\) at \(A\) and all other common clients at \(B\). Since \(C=x+z+\delta+d\), its loads are

\[
L_A=r-z-d,\qquad L_B=r-x.
\]

In fact \(L_B\ge q\). Otherwise \(x+\delta>a\), and hence \(C>a+2d\). Together with \(C<qr\) this gives

\[
d<\frac{q^3-q\delta}{2}.
\]

But \(d>U>q^4-a\delta\); combining these inequalities forces \(\delta>a\), contrary to \(\delta<a/2\). Also, since \(x\ge z\) and \(\delta\ge d\),

\[
2(z+d)\le C<qr,
\qquad L_A>(1-q/2)r>qr.
\]

Repair this pure profile by successive strict improvements. If its initial gap \(G=x-z-d\) is nonnegative, both loads are at least \(q\), and every improvement leaves both new loads strictly between the previous two loads, so both quotas remain satisfied.

If \(G<0\), its magnitude is at most \(d\). The clients \(x,z\) and all large clients have weights at least \(d\), and cannot improve; only the small clients, each below \(d/3\), can move. Before the first sign reversal \(L_A\) increases and \(L_B\ge V/2>q\). At the first reversal the positive gap is smaller than the moved weight, hence below

\[
d/3\le\delta/3<a/6<q^4/2<U.
\]

At this point both loads exceed \(q\), and remain so during all later improvements. The finite potential ensures termination at a pure Nash equilibrium meeting the two load bounds.


\subsection{A bad negative gap cannot persist}\label{a-bad-negative-gap-cannot-persist}

Suppose instead the constructed equilibrium has gap \(-e<-Z\). Every pure \(B\) client has weight at least \(z\ge e\). If any exists, the two-mixer balance \(P_A-P_B=e-\delta\) gives

\[
C=x+z+2P_B+e-\delta\ge5e-\delta,
\]

which Appendix~\ref{equilibrium-coordinates-and-three-mixers} rules out. Thus all remaining clients are at \(A\), with total \(T=e-\delta\le e\). If some has weight \(v\ge e/3\), mix it with \(x,z\), giving gap \(-(e-v)/2\). This is a valid designated-three-client equilibrium and meets the two load bounds. We may therefore assume every remaining weight is below \(e/3\).

Since \(C=x+z+e-\delta\ge2z+e-\delta\) and \(e>Z\),

\[
2z<C-Z+\delta=2C-2a+2a\delta<2q^3r.
\]

Consequently \(e\le z<q^3r\). Moreover

\[
U>q^4-a\delta>q^3r/3>e/3.
\]

For the middle inequality, multiply by three and use

\[
3q^4-q^3-(3a-q^3)\delta>
3q^4-q^3-(3a-q^3)a/2=(5-8q)/2>0.
\]

Change to the pure profile with \(x\) at \(B\), and every other common client at \(A\). Its loads are

\[
L_A=1-x\ge1-a=q,\qquad L_B=1-z-e.
\]

If \(L_B\ge q\), successive pure improvements preserve both lower bounds. Otherwise the gap \(G=z+e-x\) is positive and satisfies

\[
U<G\le e.
\]

Clients \(x,z\), having weights at least \(e\), cannot improve. Only the remaining tiny clients can move. While the gap is positive, \(L_A\ge V/2>q\), and \(L_B\) increases. At the first negative sign reversal the new absolute gap is below the weight moved, hence below \(e/3<U\le Z\); both loads then exceed \(q\), and all later improvements are safe.

If the process terminates without reaching \(L_B\ge q\), its positive gap is still larger than \(U>e/3\). No tiny client can remain at \(A\), because every one would have a strict improvement. Thus all tiny clients must be at \(B\), and \(L_B=r-z\). But

\[
\delta+z<q^3+(1-q^3)\delta
=q^3+2a\delta<q^3+a^2=a,
\]

so \(L_B=r-z>q\), a contradiction. The final equilibrium therefore meets both targets.

This proves the lemma for arbitrary finite numbers of shared clients.

\subsection{Termination of pure improvements}

If a shared client of weight \(w\) strictly improves by leaving a facility of load \(H\) for one of load \(L<H\), then \(w<H-L\). The new loads \(H-w,L+w\) both lie strictly between \(L,H\). The absolute load gap strictly decreases; if its sign reverses, the new absolute gap is less than \(w\). The potential \(L_A^2+L_B^2\) strictly decreases, so any sequence of strict pure improvements terminates in this finite game. All constructions and repairs above therefore use ordinary independent mixed equilibria or pure equilibria.



\section{A polynomial constructor for the strong chord}
\label{app:polynomial-chord}

This appendix supplies the algorithmic step used in menu C. The objective
is to obtain the exchange inequalities used in Appendix~\ref{largest-client-and-the-maximum-square-construction}. Computing a
global maximum of a convex quadratic function is unnecessary.

\subsection{An exchange algorithm on a box slice}
Let $v_1,\ldots,v_m>0$, let $|t|\le\sum_i v_i$, and define
\[
 P=\{z\in\mathbb R^m:-v_i\le z_i\le v_i,\ \sum_i z_i=t\},
 \qquad F(z)=\sum_i z_i^2.
\]
A vertex has at most one coordinate strictly between its two bounds. To
obtain an initial vertex, begin at $z_i=-v_i$ and increase coordinates in
index order until their total increase is $t+\sum_i v_i$. All but possibly
one coordinate end at a bound. If all coordinates are at bounds, stop.
Otherwise let $j$ be the unique interior coordinate, write $v_j=z$ and
$z_j=d$, and call $j$ the pivot.

For each $i\ne j$ write $z_i=sv$, where $s\in\{-1,1\}$ and $v=v_i$.
Holding the sum $b=d+sv$ fixed, the feasible pair segment is
\[
 z_j=u,\qquad z_i=b-u,\qquad
 u\in[\max(-z,b-v),\ \min(z,b+v)].
\]
The current $d$ is one endpoint. Inspect the other endpoint. If it strictly
increases the two squares, move there and restart the scan in index order.
If no pair endpoint strictly improves, return the current vertex. At each
move at least one of the two changed coordinates is at a bound; all other
coordinates were already at bounds. Thus feasibility and the vertex
invariant are preserved.

\begin{theorem}\label{thm:exchange}
This algorithm terminates after $O((m+1)^3)$ moves and
$O((m+1)^4)$ rational arithmetic operations. Its output is a vertex at
which no two-coordinate endpoint move increases $F$.
\end{theorem}
\begin{proof}
If a nonterminal move changes the pivot from weight $z$ to weight $v$,
the old pivot reaches the bound $sz$. The gain is
\[
 2(z-sd)(z-v).
\]
Since $|d|<z$, a strict gain implies $v<z$. Every pivot change strictly
decreases its weight, so there are at most $m-1$ changes.

If the pivot does not change, coordinate $i$ flips from $sv$ to $-sv$,
$d'=d+2sv$, and $|d'|>|d|$. Fix a maximal phase with one pivot.
When such a flip reverses the nonzero sign of $d$, it requires $v>|d|$
and leaves $|d'|=2v-|d|>v$. The weights of successive sign-reversing
flips therefore strictly increase. There are at most $m-1$ reversals in
one phase. Between reversals every nonpivot coordinate flips at most
once: a nonreversing improving flip must have $s$ aligned with the sign
of $d$, and afterwards that coordinate is opposed to that sign. A first
flip from $d=0$ merely establishes a sign; a strict gain never returns
to zero. Thus each phase contains $O((m+1)^2)$ flips.
There are at most $m$ phases and at most one final all-bound move.
Scanning at most $m-1$ pairs for each move proves the claimed operation
bound. The estimate is independent of weight magnitudes and separations.
\end{proof}

\subsection{The exact inequalities supplied by the output}
Suppose the returned vertex still has pivot weight $z$ and value $d$.
On each pair segment, $F$ is convex. Since the opposite endpoint has no
larger value, no interior point has larger value either. Its directional
derivative at the current endpoint therefore gives
\begin{equation}
 z_i=+v\Rightarrow d\le v,\qquad
 z_i=-v\Rightarrow d\ge-v.                         \label{eq:old-B-1}
\end{equation}
For $d>0$ the endpoint comparisons also imply
\begin{equation}
 z_i=+v\Rightarrow v\ge z,\qquad
 z_i=-v,\ v<z\Rightarrow v\le d.                  \label{eq:old-B-2}
\end{equation}
Indeed, for a positive coordinate, if $2v\ge z-d$, moving the pivot to
$z$ has gain $2(z-d)(z-v)$; otherwise flipping the coordinate has gain
$4v(d+v)>0$. Either alternative excludes $v<z$. For a negative coordinate
with $v<z$, if $2v\ge z+d$, moving the pivot to $-z$ has the forbidden
positive gain $2(z+d)(z-v)$. Otherwise its full flip has gain
$4v(v-d)$, forcing $v\le d$. For $d<0$ the signs reverse: pure negative
coordinates have weight at least $z$, and pure positive coordinates
below $z$ have weight at most $-d$.

These are exactly the local derivative and four endpoint comparisons in
Appendix~\ref{largest-client-and-the-maximum-square-construction}. That appendix proves the initial customer equilibrium using
\eqref{eq:old-B-1}. Its positive-gap branch uses the two statements of \eqref{eq:old-B-2}; its
negative-gap branch uses their sign reversal. The later mass inequalities,
explicit mixer choices and pure repairs never compare quadratic objectives
again. Thus every remaining argument in Appendix~\ref{app:cross-chord} is valid for this
polynomially obtained vertex.

The output need not maximize $F$ globally. For weights $(14,13,19)$ and
target $-10$, the vertex $(14,-5,-19)$ has no improving pair endpoint,
but its value $582$ is below the value $654$ at $(-14,-13,17)$.
This distinction is why only the required exchange inequalities are claimed.

\subsection{Completing and timing the strong-chord construction}
Normalize the larger reach to one, put $\delta=1-r$, choose a largest
common weight $x$, and write $S$ for all other common weight. If
$S\ge\delta$, apply Theorem~\ref{thm:exchange} to the other coordinates
with target $-\delta$. If its output is all-bound, give $x$ coordinate
zero. Otherwise give $x$ the pivot coordinate $d$. The coordinate of
customer $i$ is $c_i=w_i(2p_i(A)-1)$, and the gap identity is
$\Delta=\delta+\sum_i c_i$. Therefore the resulting initial profile is
exactly the equilibrium of Appendix~\ref{largest-client-and-the-maximum-square-construction}, with one or two designated
mixers. If it meets the two quotas, return it.

If a quota fails, execute the following ordered branch table from Appendices~\ref{a-bad-positive-gap-cannot-persist}--\ref{a-bad-negative-gap-cannot-persist}. For a positive gap \(d>U\), first test for a pure \(A\)-side client; the three-mixer bound rules this case out. Otherwise scan for a nonlarge client of weight in \([d/3,d]\); if one exists, use it as the third mixer. If none exists, sort the remaining weights, take the prefix whose running sum first reaches \(d\), and apply the large-client split and the final pure repair in the order stated in the appendix. For a negative gap \(-e<-Z\), test for a pure \(B\)-side client, then scan for a remaining client of weight at least \(e/3\), and otherwise use the symmetric pure seed and repair. Every test is an exact comparison of rational sums; every prefix and sort is linear or \(O(k\log k)\), and a pure branch always moves a largest strictly improving client. The quota-preservation arguments in Appendix~\ref{app:cross-chord} apply to each returned branch.
If $S<\delta$, fix $x$ at the lower-reach location and perform the same
largest-improvement rule on the other customers, as in Appendix~\ref{largest-client-and-the-maximum-square-construction}.
The fixed customer's stability is proved there and is preserved throughout.

For completeness, this pure improvement rule moves at most $k$ customers
in a two-facility game with $k$ movable common customers, even with private
base loads. An improving weight $w$ changes absolute gap $g$ to
$|g-2w|<g$. Before a sign reversal, candidates on the higher side can only
disappear, so chosen weights are nonincreasing. At a reversal the new gap
is below the last moved weight; every earlier moved customer is therefore
permanently stable, as all future gaps are smaller. Any newly movable
customer has weight below the new gap, hence below the preceding chosen
weight. Inductively no customer moves twice. A direct scan for each move
costs $O((k+1)^2)$ rational operations. Fixed customers are simply absorbed
into the private base loads in this argument.

The exchange algorithm consequently dominates the running time:
$O((k+1)^4)$ rational operations per strong chord. Every exchange coordinate
is an endpoint or the slice target minus a signed sum of input weights.
Subsequent mixed gaps divide signed sums by integers at most $k$, and
probabilities divide by the corresponding input weight. The encoding
lengths remain polynomial. The exact golden-ratio comparisons described
in Section~\ref{sec:algorithm} avoid numerical tolerances.

The elementary two-facility improvement analysis above is also consistent
with the maximum-weight improvement rule studied for load balancing by
Even-Dar, Kesselman, and Mansour~\cite{evendar2007}. The box-slice exchange
bound is supplied in full here; no efficient global convex-maximization
claim is used.


\bibliographystyle{plain}
\bibliography{references}
\end{document}
